\section{Estimating the Environmental Impact of Training our Models} 

\begin{table}[t]
  \centering
  \begin{tabular}{@{}l cccccc@{}}
    \toprule
    Model to      & \multirow{2}{*}{GPU Type} & GPU Power &  \multirow{2}{*}{GPU-hours} & \multirow{2}{*}{PUE} & Total power  & Carbon emitted\\
    Reproduce      &  & consumption & &      &  consumption & (tCO$_2$eq)\\
    \midrule
    \OURS-g     & A100-40GB & 400W & 22,016 & 1.1 & 9.7 MWh  & 3.7 \\       
    \bottomrule
  \end{tabular}
  \caption{
    \textbf{Carbon footprint of reproducing \OURS.}
    We report the potential carbon emission of reproducing \OURS-g when assuming a  power consumption for the A100-40GB of 400W,  a PUE of 1.1 and carbon intensity factor of 0.385 kg CO$_2$e per KWh.
    \label{tab:carbon}
  }
\end{table}

Training foundation models consumes a significant amount of energy, resulting in carbon dioxide emissions.
\citet{patterson2021carbon} propose a methodology to report an estimation of the carbon emitted during the training of a model based on the specifics of the data center and its power grid.
This computation informs the design of the data center used for the training of models and the choice of location for data centers.
This methodology requires to know the specifics of the data center used for training, which can be complex when multiple data centers are involved over time. 
Additionally, these specifics are most often not in the control of the AI practitioner, and hence, this methodology is less helpful when practioners make technical decisions about future trainings.
Instead, in this section, we follow an alternative that reports the potential carbon emission of retraining a similar model in an average data center located in the US.
This methodology was used in previous work in natural language processing~\citep{strubell2019energy,touvron2023llama} to establish an apple-to-apple comparison between pretraining schemes.
More precisely, we fix the value of all exogenous variables, i.e., the Power Usage Effectiveness (PUE) and carbon intensity factor of a power grid to the same values as in \citet{touvron2023llama}, that is, a PUE of 1.1 and the carbon intensity factor to the US average of 0.385 kg CO$_2$eq/KWh.
We use the same formula as in~\cite{patterson2021carbon} to estimate the potential energy consumption and the carbon emission.
For the power consumption of an A100-80GB, we take the thermal design power for NVLink systems, which is 400W.
We report the potential carbon emission of retraining a \OURS{} ViT-g in Table~\ref{tab:carbon}.
For comparison, retraining an OpenCLIP ViT-L or OpenCLIP ViT-G would require 22.4 MWh and 118.9 MWh, respectively, if run in the same data center. 
This is 10$\times$ more carbon emission.
Note that this comparison is not fair to them, since they also train a text encoder in parallel, and we thus do not report them in the table.
However, it gives a reasonable \rthree{guideline} for those who are interested in training only visual features: in this context, training a self-supervised model is preferable in terms of carbon emission.
Training a text-guided model still makes sense when planning to reuse the text encoder.


\paragraph{Carbon footprint of the whole project.} 
Additionally, we estimate the footprint of the whole project to be between $0.5$k and $1$k tCO$_2$eq using the same grid as presented above \footnote{\rthree{For context, a full Boeing 777 return flight between London and New York corresponds to approximately 560 tCO$_2$eq.}}. 
This carbon footprint represents in the order of $200$k GPU-days.
The primary sources of emissions are the self-supervised pre-trainings of the models.
For example, a single pre-training of a ViT-g model (22k GPU-hours) emits 3.7 tons of CO$_2$eq, while a finetuning on ImageNet-1k (1k GPU-hours) emits 0.2 tons.
This estimate only considers the GPUs' electricity consumption and ignores other emissions, such as their manufacturing and disposal.

